\scanpage{43}
\mspage{39}
% archive sheet number: 42

\[
\unsure{\Sigma\mathrm{cons}_{\mathcal A}|\Sigma_{\mathrm{cons}_{\mathcal A}}}
\]
Elle n'est \deleted{\unsure{en}} en général stable ni par $\mathrm{Sup}$ ni par $\mathrm{Inf}$ ni par $C$. Mais il y a un axiome (``axiome des supports'') qui tient lieu d'une telle compatibilité, sur laquelle nous allons avoir à revenir --- et à vérifier \inserted{encore} cet axiome, en faisant intervenir cette fois aussi la structure $\unsure{\overset{\circ}{\ll}|\overset{\circ}{\Subset}}$.

\deleted{Pour l'instant, je suis \unsure{resté}} \inserted{Du point de vue notation} pour $X\in\mathcal A$, je dénote par
\[
-\underset{\substack{\Vert\\\inserted{\mathrm{supp}^{\circ}(X)}}}{X^{\circ}}\in\Sigma\mathrm{cons}_{\mathcal A}
\]
l'élément $\mathrm{supp}^{\circ}X$, qui peut être décrit comme le plus petit $S\in\Sigma_{\mathcal A}$ (donc aussi de $\Sigma\mathrm{cons}_{\mathcal A}$) tel que $X\in S$. C'est une explicitation possible de l'intuition de \unsure{la|l'} \uline{\unsure{Art|Aut} ouvert} associée à un Art donné $X$.
